\documentclass[11pt]{article}
\usepackage[a4paper,margin=26mm]{geometry}
\usepackage{amsmath,amssymb}
\usepackage[hidelinks]{hyperref}
\setlength{\emergencystretch}{3em}
\title{Conjecture 00000001670: incompatible numerical characterizations}
\author{lizaixi01, with OpenAI Codex assistance}
\date{23 October 2026}
\begin{document}
\maketitle
\section{The necessary numerical clauses}
The English source assigns the same real constant $P_3$ both values
\[
 P_3=0.918=\frac{459}{500},\qquad
 P_3=\frac{2}{\pi}\arctan\sqrt2.
\]
The Chinese source also prints these two numerical characterizations.
They cannot hold simultaneously. We disprove this necessary part of the
conjecture, hence the stated conjunction, without assuming or formalizing
a random-graph limit. We do not claim to determine a genuine max-cut
constant, or to disprove the separate asymptotic or concentration clauses.

Even interpreting $0.918$ as a three-decimal approximation does not repair
the contradiction. We allow the inclusive interval
\[
 \left|P_3-\frac{459}{500}\right|\le\frac1{2000},
 \qquad\text{equivalently}\qquad
 0.9175\le P_3\le0.9185.
\]
This includes either possible tie at a rounding endpoint.

\section{A strict analytic bound}
Set $c=(2/\pi)\arctan\sqrt2$. The ordinary nonnegative real square
root is strictly increasing on positive inputs, so $\sqrt2<\sqrt3$.
The principal real inverse tangent is strictly increasing and satisfies
$\arctan\sqrt3=\pi/3$. Since $\pi>0$, multiplication by $2/\pi$
preserves the strict inequality and gives
\[
 c<\frac2\pi\arctan\sqrt3
   =\frac2\pi\frac\pi3=\frac23.
\]
Every function and constant in this chain is Mathlib's actual real
analytic object in the Lean proof, not an assumed numerical relation.

The rational comparisons now leave a large gap:
\[
 \frac{459}{500}-c>
 \frac{459}{500}-\frac23=\frac{377}{1500}>\frac14.
\]
Under the rounding interpretation, $P_3\ge367/400$ and therefore
\[
 P_3-c>\frac{367}{400}-\frac23
   =\frac{301}{1200}>\frac14.
\]
Thus the exact closed form differs by more than $1/4$ even from every
value in the permitted rounding interval. No numerical approximation
to $\pi$, $\sqrt2$, or $\arctan\sqrt2$ is needed.

\newpage
\section{Formal scope and previous submission}
The formal definition \texttt{closedForm} is literally
\texttt{(2 / Real.pi) * Real.arctan (Real.sqrt 2)} in $\mathbb R$.
\texttt{Rounded0918} is the absolute-error interval above.
The file proves the square-root comparison, inverse-tangent comparison,
strict $2/3$ bound, and the two strict error bounds before proving
\[
 \forall P\in\mathbb R,\quad
 \neg\bigl(P=459/500\ \wedge\ P=c\bigr)
\]
and
\[
 \forall P\in\mathbb R,\quad
 \neg\bigl(|P-459/500|\le1/2000\ \wedge\ P=c\bigr).
\]
The final two theorems also negate existence of any sequence
$P:\mathbb N\to\mathbb R$ satisfying the respective requirements at
index $3$. These universally quantified statements apply to any
proposed real $P_3$; no special definition of a graph parameter is
inserted to make the contradiction hold.

\href{https://github.com/The-Last-Math-Competition/The-Last-Math-Competition/pull/189}
{Closed PR 189}, by orionsheep, already presented the same elementary
analytic comparison in prose. We acknowledge that prior observation;
this submission does not claim it as a new discovery. The
\href{https://github.com/The-Last-Math-Competition/The-Last-Math-Competition/pull/189\#issuecomment-5965633103}
{maintainer's review} rejected its Lean component because it proved
only natural-number comparisons and reflexive equalities. The
load-bearing real $\pi$, square root, inverse tangent, and $P_3$
contradiction were absent from Lean. The present proof formalizes that
entire chain and additionally covers the rounding interpretation.
It relies on the existing, proved Mathlib analytic theorems rather
than treating the analytic steps as assumptions.

\section{Background and verification}
Dembo, Montanari and Sen,
\href{https://arxiv.org/abs/1503.03923}{\emph{Extremal Cuts of Sparse
Random Graphs}}, study a large-degree asymptotic governed by a constant
$P_*$. Their result provides background for the source's terminology.
We do not identify $P_*$ with the source's $P_3$, use a reported decimal
for $P_*$ as a premise, or claim to formalize their probabilistic result.
The disproof here is solely the inconsistency of the two printed
necessary numerical characterizations.

Lean 4.33.0 and Mathlib commit
\texttt{db584cd6d46c92f209a44c0f1c829460d327499d} are pinned.
The package provides a new Main build, ten transitive axiom audits,
and full official \texttt{leanchecker --fresh Main} replay, with
source-bound commands, exit codes and logs. No \texttt{sorry},
\texttt{native\_decide} or additional axiom is used.
The LaTeX source and rendered PDF are included. OpenAI Codex assisted;
a separate supervising Codex agent reviewed the formal chain and report.
This is AI review, not human certification. Organizer acceptance is pending.
\end{document}
